% Part: first-order-logic
% Chapter: syntax-and-semantics
% Section: Structures

\documentclass[../../../include/open-logic-section]{subfiles}

\begin{document}

\olfileid{fol}{syn}{str}

\olsection{మొదటిస్థాయి విధేయ భాషల కోసం \printtoken{P}{structure}}

\begin{explain}
మొదటిస్థాయి విధేయ భాషలకు వాటంతట వాటికి \emph{అర్థనిర్దేశం ఉండదు:}
వాటిలోని \tetoken{స్థిరసంకేతాలు}{!!{constant}s}, \tetoken{ప్రమేయ సంకేతాలు}{!!{function}s}, \tetoken{విధేయ సంకేతాలు}{!!{predicate}s} నిర్దిష్టమైన
అర్థం జోడించబడి ఉండదు. ఒక \article{structure} \emph{\tetoken{నిర్మాణం}{!!{structure}}}ను
నిర్దేశించడం ద్వారా వాటికి అర్థాలు ఇస్తాం. అది \emph{వ్యక్తి క్షేత్రం}ను,
అంటే \tetoken{స్థిరసంకేతాలు}{!!{constant}s} సూచించే, \tetoken{ప్రమేయ సంకేతాలు}{!!{function}s} పనిచేసే, పరిమాణకాలు
పరిధి కలిగి ఉండే వస్తువులను నిర్దేశిస్తుంది. అదనంగా, ఏ \tetoken{స్థిరసంకేతాలు}{!!{constant}s}
ఏ వస్తువులను సూచిస్తాయో, ఒక \tetoken{ప్రమేయ సంకేతం}{!!a{function}} వస్తువులను వస్తువులకు ఎలా
పంపుతుందో, \tetoken{విధేయ సంకేతాలు}{!!{predicate}s} ఏ వస్తువులకు వర్తిస్తాయో అది నిర్దేశిస్తుంది.
తర్కంలో పర్యవసానం, చెల్లుబాటుతనం, సంతృప్తిపరచదగినతనం వంటి
\emph{అర్థపర} భావనలకు \tetoken{నిర్మాణాలు}{!!^{structure}s} ఆధారం. సాహిత్యంలో వీటిని
``నిర్మాణాలు,'' ``అర్థనిర్దేశాలు,'' లేదా ``నమూనాలు'' అని వివిధంగా
పిలుస్తారు.
\end{explain}

\begin{defn}[\tetoken{నిర్మాణాలు}{!!^{structure}s}]
మొదటిస్థాయి విధేయ తర్క భాష~$\Lang{L}$కు ఒక \Article{structure}
\emph{\tetoken{నిర్మాణం}{!!{structure}}}~$\Struct M$ కింది అంశాలతో కూడి ఉంటుంది:
\begin{enumerate}
\item \emph{వ్యక్తి క్షేత్రం:} ఖాళీ కాని సమితి, $\Domain M$
\item \emph{\tetoken{స్థిరసంకేతాల}{!!{constant}s} అర్థనిర్దేశం:} $\Lang{L}$లోని ప్రతి
  \tetoken{స్థిరసంకేతం}{!!{constant}}~$c$కు ఒక \tetoken{మూలకం}{!!a{element}} $\Assign{c}{M} \in \Domain M$
\item \emph{\tetoken{విధేయ సంకేతాల}{!!{predicate}s} అర్థనిర్దేశం:} $\Lang{L}$లోని ప్రతి
  $n$-స్థాన \tetoken{విధేయ సంకేతం}{!!{predicate}}~$R$కు ($\eq$ తప్ప), ఒక $n$-స్థాన సంబంధం
  $\Assign{R}{M} \subseteq \Domain{M}^n$
\item \emph{\tetoken{ప్రమేయ సంకేతాల}{!!{function}s} అర్థనిర్దేశం:} $\Lang{L}$లోని ప్రతి
  $n$-స్థాన \tetoken{ప్రమేయ సంకేతం}{!!{function}}~$f$కు, ఒక $n$-స్థాన ప్రమేయం
  $\Assign{f}{M} \colon \Domain{M}^n \to \Domain{M}$
\end{enumerate}
\end{defn}

\begin{ex}
అంకగణిత భాషకు ఒక \tetoken{నిర్మాణం}{!!^a{structure}}~$\Struct M$, ఒక సమితి, ఆ సమితి
$\Domain M$లోని ఒక మూలకం $\Assign{\Obj 0}{M}$ (ఇది
\tetoken{స్థిరసంకేతం}{!!{constant}}~$\Obj 0$ యొక్క అర్థనిర్దేశం), ఒక ఏకస్థాన ప్రమేయం
$\Assign{\Obj \prime}{M} \colon \Domain{M} \to \Domain M$, రెండు
ద్విస్థాన ప్రమేయాలు $\Assign{\Obj +}{M}$, $\Assign{\Obj
  \times}{M}$ (రెండూ $\Domain M^2 \to \Domain M$), ఇంకా ఒక ద్విస్థాన
సంబంధం $\Assign{\Obj <}{M} \subseteq \Domain{M}^2$తో కూడి ఉంటుంది.

అటువంటి నిర్మాణానికి స్పష్టమైన ఉదాహరణ ఇది:
\begin{enumerate}
\item $\Domain N = \Nat$
\item $\Assign{\Obj 0}{N} = 0$
\item అన్ని $n \in \Nat$లకు $\Assign{\Obj \prime}{N}(n) = n + 1$
\item అన్ని $n, m \in \Nat$లకు $\Assign{\Obj +}{N}(n, m) = n + m$
\item అన్ని $n, m \in \Nat$లకు $\Assign{\Obj \times}{N}(n, m) = n\cdot m$
\item $\Assign{\Obj <}{N} = \Setabs{\tuple{n, m}}{n \in \Nat, m \in
  \Nat, n < m}$
\end{enumerate}
ఈ విధంగా నిర్వచించిన $\Lang L_A$ యొక్క నిర్మాణం~$\Struct N$ను
\emph{అంకగణిత ప్రామాణిక నమూనా} అంటారు; ఎందుకంటే అది~$\Lang L_A$లోని
తార్కికేతర స్థిరసంకేతాలకు మనం సహజంగా ఆశించే అర్థాన్నే ఇస్తుంది.

అయితే~$\Lang L_A$కు మరెన్నో \tetoken{నిర్మాణాలు}{!!{structure}s} సాధ్యమే. ఉదాహరణకు,
వ్యక్తి క్షేత్రంగా సహజ సంఖ్యల సమితి~$\Nat$కు బదులు పూర్ణాంకాల
సమితి~$\Int$ను తీసుకొని, $\Obj 0$, $\Obj \prime$, $\Obj +$,
$\Obj \times$, $\Obj <$లకు తగిన అర్థనిర్దేశాలు ఇవ్వవచ్చు. సంఖ్యలతో
ఏమాత్రమూ సంబంధం లేని నిర్మాణాలను కూడా~$\Lang L_A$కు నిర్వచించవచ్చు.
\end{ex}

\begin{ex}
సమితి సిద్ధాంత భాష~$\Lang L_Z$కు ఒక నిర్మాణం~$\Struct M$లో కేవలం
ఒక సమితి, ఒక్క ద్విస్థాన సంబంధం అవసరం.\sourcecorrection{OLTEFOLSEM-001}{మూలంలోని ``single-two place relation''ను సందర్భానుసారంగా ``single two-place relation''గా చదివాం.} అందువల్ల సాంకేతికంగా,
ఉదాహరణకు మనుషుల సమితితో పాటు ``$x$, ~$y$ కంటే పెద్దవయస్కుడు'' అనే
సంబంధాన్ని, లేదా $n, m \in \Nat$లకు $n \ge m$ అనే సంబంధంతో పాటు
$\Nat$ను, $\Lang L_Z$కు ఒక \tetoken{నిర్మాణంగా}{!!a{structure}} ఉపయోగించవచ్చు.

$\Lang L_Z$కు ప్రత్యేకంగా ఆసక్తికరమైన \tetoken{నిర్మాణంలో}{!!{structure}}
వ్యక్తి క్షేత్రంలోని \tetoken{మూలకాలు}{!!{element}s} నిజంగానే సమితులు; $\Obj \in$ యొక్క
అర్థనిర్దేశం నిజంగానే ``$x$, ~$y$లోని ఒక \tetoken{మూలకం}{!!a{element}}'' అనే సంబంధం.
అది \emph{వంశపారంపర్యంగా పరిమిత సమితుల} \tetoken{నిర్మాణం}{!!{structure}}~$\Struct{{HF}}$:
\begin{enumerate}
\item $\Domain{{{HF}}} = \emptyset \cup \Pow{\emptyset} \cup
  \Pow{\Pow{\emptyset}} \cup \Pow{\Pow{\Pow{\emptyset}}} \cup \dots$;
\item $\Assign{\Obj \in}{{{HF}}} = \Setabs{\tuple{x, y}}{x, y \in
  \Domain{{{HF}}}, x \in y}$.
\end{enumerate}
\end{ex}

\begin{digress}
ఏది ఒక \tetoken{నిర్మాణంగా}{!!a{structure}} పరిగణించబడుతుందో మనం చేసే నిర్దేశాలు మన
తర్కంపై ప్రభావం చూపుతాయి. ఉదాహరణకు, ఖాళీ వ్యక్తి క్షేత్రాలను అనుమతించకపోవడం
వల్ల, పరిమాణీకృత వాక్యాల సంతృప్తి (లేదా సత్యం) గురించి సాధారణంగా
ఇచ్చే వివరణ ప్రకారం, $\lexists[x][(!A(x) \lor \lnot !A(x))]$
చెల్లుబాటు అవుతుంది---అంటే అది ఒక తార్కిక సత్యం. అంతేకాక, అన్ని
\tetoken{స్థిరసంకేతాలు}{!!{constant}s} తప్పనిసరిగా వ్యక్తి క్షేత్రంలోని ఒక వస్తువును సూచించాలనే
నిబంధన వల్ల, అస్తిత్వ సాధారణీకరణ ఒక నిర్దుష్ట నిగమన నమూనా అవుతుంది:
$!A(a)$, కాబట్టి $\lexists[x][!A(x)]$. పేర్లు వ్యక్తి క్షేత్రం బయటివాటిని
సూచించడానికి లేదా దేనినీ సూచించకపోవడానికి అనుమతిస్తే, మనం
\emph{స్వేచ్ఛా తర్కం} వైపు వెళ్తాం. అందులో అస్తిత్వ సాధారణీకరణకు ఒక
అదనపు పూర్వాధారం అవసరం: $!A(a)$ మరియు $\lexists[x][\eq[x][a]]$,
కాబట్టి $\lexists[x][!A(x)]$.
\end{digress}

\end{document}
